\documentclass[a4paper,11pt]{article}

% ==================================================
% 1. Packages
% ==================================================

\usepackage{kotex}
\usepackage{amsmath}
\usepackage{mynotes}

% ==================================================
% 2. Document
% ==================================================

\begin{document}

% ==================================================
% 3. Cover Page
% ==================================================

\begin{titlepage}
    \centering
    \vspace*{\fill}

    {\Huge\bfseries Title\par}

    \vspace{1cm}

    {\Large Author\par}

    \vspace{0.5cm}

    {\large \today\par}

    \vspace*{\fill}
\end{titlepage}

% ==================================================
% 4. Lecture Notes
% ==================================================

% --------------------------------------------------
% 4-1. Section / Subsection
% --------------------------------------------------

\section{Section}

\subsection{Subsection}

\subsubsection{Subsubsection}

일반적인 본문 내용을 작성한다.


% --------------------------------------------------
% 4-2. Paragraph Spacing
% --------------------------------------------------

첫 번째 문단이다.

\dhknp

두 번째 문단이다.


% --------------------------------------------------
% 4-3. Horizontal Spacing
% --------------------------------------------------

\dhkti 1cm 가로 공백을 적용한 문장이다.


% --------------------------------------------------
% 4-4. Short Note Box
% --------------------------------------------------

\dhknt{짧은 회색 강조 박스 내용}


% --------------------------------------------------
% 4-5. Multiline Note Box
% --------------------------------------------------

\begin{dhknote}
여러 줄 회색 강조 박스 내용이다.

긴 설명이나 여러 문단을 작성할 수 있다.
\end{dhknote}


% --------------------------------------------------
% 4-6. Inline Code
% --------------------------------------------------

Python의 \texttt{sys.path}는 모듈 검색 경로를 나타낸다.


% --------------------------------------------------
% 4-7. Python Code Block
% --------------------------------------------------

\begin{dhkcode}
import sys

print(sys.version)
print(sys.executable)
print(sys.path)
\end{dhkcode}


% --------------------------------------------------
% 4-8. Tree Diagram
% --------------------------------------------------

\begin{forest}
dhktree
[Python
    [표준 라이브러리
        [sys]
        [os]
        [pathlib]
    ]
    [외부 라이브러리
        [numpy]
        [pandas]
    ]
]
\end{forest}


% --------------------------------------------------
% 4-9. Text Formatting
% --------------------------------------------------

일반 텍스트

\textbf{굵은 텍스트}

\textit{기울임 텍스트}

\texttt{고정폭 텍스트}


% --------------------------------------------------
% 4-10. Center Alignment
% --------------------------------------------------

\begin{center}
가운데 정렬된 텍스트
\end{center}


% --------------------------------------------------
% 4-11. Page Break
% --------------------------------------------------

\newpage

\section{Next Section}

새로운 페이지에서 시작하는 내용이다.


% ==================================================
% 5. End Document
% ==================================================

\end{document}